
\subsection{Distinguishing neural representations across mouse genotypes}

In every dataset so far, the systems we compared formed a continuum: cortical areas lie along a gradient tracking hierarchy and tuning properties, not in discrete groups. In what follows, we show that shape metrics can find clusters in subject space in a cohort of animals in which the same visual discrimination task (\textbf{Fig.} \ref{fig:ibl}a) was performed by
wild-type mice and by three mouse lines carrying mutations associated with autism
spectrum disorder \hl{(\textit{C57BL/6}, $n=10$; \textit{Fmr1}, $n$ = 9; \textit{Cntnap2}, $n$ = 11 and \textit{Shank3}, $n$ = 9)} released by \cite{noel2025}.

% That study characterized these genotypes through the lens of prior encoding. Stimulus processing, by contrast, was reported to be unaffected: psychometric performance did not differ across genotypes, and inspection of example rasters and peri-stimulus time histograms showed the same standard response features in every mouse line.
% That comparison, however, was made at the level of single neurons. Population geometry is a different level of description \parencite{Kriegeskorte2021}, and whether it also is indistinguishable across genotypes has not been asked. Moreover, an important question remains open: the original analysis pools units within a genotype, so no individual animal is ever compared to another individual animal, arguably the natural unit for asking whether the three lines differ from controls.

For each animal, we assembled a matrix of stimulus-evoked encoding kernels across neurons with 106 time bins (as in \cite{noel2025}), and retained animals with at least n = 50 units, \hl{leaving $n=9$ control animals and $n=9$, $n=6$, and $n=9$ for \textit{Fmr1}, \textit{Cntnap2}, and \textit{Shank3}, respectively}. We then computed a $K \times K$ matrix of pairwise shape distances across all animals using Procrustes and pooling genotypes, so that each animal occupies a single position in a common space. Hierarchical clustering of this matrix suggested that wild-type animals are clustered together (\textbf{Fig.} \ref{fig:ibl}b). We quantify this with two complementary analyses, clustering and decoding in subject space. First, the best silhouette in subject space exceeded the null ($z = +4.78$, $p < 0.01$; \textbf{Fig.}~\ref{fig:ibl}c, \textit{left}), so the animals fall into discrete groups rather than forming one continuous cloud. Repeating the same test on the mutant animals alone removed the effect ($z = +0.86$, $p = 0.16$; \textbf{Fig.}~\ref{fig:ibl}c, \textit{right}), indicating that what separates animals in this space is wild type from mutant, not one mutant line from another. Second, we decoded genotype directly from the distance matrix: each held-out mouse was assigned to whichever group the remaining animals were closer to on average in shape space, and each of four one-versus-rest contrasts was scored by balanced accuracy against $n=$10000 label permutations. Wild-type animals were identified well above chance (accuracy =  $72\%$, $p = 0.010$, \textbf{Fig.} \ref{fig:ibl}h), whereas none of the three mutant lines was distinguishable from the remaining animals (\textit{Fmr1}, $0.44$, $p = 0.76$; \textit{Cntnap2}, $0.45$, $p = 0.74$; \textit{Shank3}, $0.51$, $p = 0.49$), reinforcing that the lines depart from controls without differing from one another. This last observation is in line with the picture of convergent phenotypes suggested by \cite{noel2025}: distinct genetic perturbations may converge on a common computational anomaly, reflected here in similar population geometry.


\begin{figure}
    \centering
     \includegraphics[width=0.99\textwidth]{content/figures/ibl.pdf}

    % \caption{   \textbf{Behavioral strategy and genetic background are reflected in neural geometry.}
    % \textbf{a.} IBL dataset schematics.
    % \textbf{b.} Psychometric curves all animals (n=56) in gray and grand average in blue.
    % \textbf{c.} Matrix of neural and behavioral dissimilarities between all animals. Note that in all the analyses, independent data was used to compute behavioral and neural distances. Here shown for one data split.  \textbf{d.} Left, correlation between behavioral and neural distances of one fold. Right, correlation values of n=100 repetitions of the data split and similar analyses for shuffled data. Both done using neural data around the response ([0-0.4 s]). \textbf{e.} The correlation between neural and behavioral distances throughout the trial and aligned to stimulus or response. Red, blue and black ticks on the top mark the time-points where stimulus, response or stimulus vs response was significantly different than 0 (bootstrap t-test), respectively. \textbf{f.} Alignment between neural representations allows for 0-shot learning of psychometric curves using the 2-closest neighbors. Behavioral correlation between the first (2 closest neighbors) and last point (2 furthest neighbors) is significantly different ($t=2.2, p=0.037$). }
    \caption{\textbf{Behavioral strategy and genetic background are reflected in neural geometry.}
    \textbf{a.} IBL datasets II \& III. Top, task. Bottom,
    Neuropixels insertions: one location replicated across \hl{n = 56} animals
    \parencite{iblreproducibility} and across the four genotypes of \parencite{noel2025}.
    \textbf{b.} Hierarchical clustering of
    pairwise shape distances across animals; black, wild type; grey, mutants.
    \textbf{c.} Animals fall into discrete groups in shape space (left, all
    animals, \hl{n = 36}: best silhouette, red, exceeds a one-blob null, blue;
    $p < 0.01$), but mutants alone do not (right, \hl{n = 26}; $p = 0.16$).
    \textbf{d.} Genotype decoding by leave-one-out nearest-group assignment,
    scored against label shuffling ($n = 10{,}000$); only wild type exceeds
    chance ($p = 0.01$). Error bars, bootstrap s.e.m. \textbf{e.} Mean shape
    distance between the five regions recorded in all \hl{n = 17} mice; Wilcoxon
    signed-rank across mice; $^{*}p<0.05$, $^{**}p<0.01$, $^{***}p<0.001$.
    \textbf{f.} Psychometric curves for all \hl{n = 56} mice (grey, individual;
    teal, mean). \textbf{g.} Pairwise distances across the 56 mice: lower
    triangle, neural shape distance (greys); upper triangle, behavioural distance
    (purples). \textbf{h.} Neural against behavioural distance for all 1,540
    mouse pairs ($r = 0.11$, $p = 1.5\times10^{-5}$). \textbf{i.}
    Shuffle-corrected behaviour--neural correlation per region (z-scored against
    shuffled distributions across cross-validation folds; circle, mean; whiskers,
    95\% C.I.).}

    \label{fig:ibl}
\end{figure}



Thus, population geometry distinguishes genotypes that appear identical at the single-neuron level \parencite{noel2025}, mirroring what we found across regions, where geometry recovered an anatomical hierarchy that single-variable decodability had obscured. In both cases, structure invisible to single-neuron analysis becomes visible once neural systems are compared as geometric objects.


\subsection{Individual differences in behavior are mirrored in population geometry}


Individual differences in geometry need not be confined to a single region. We next asked, in a final analyses, whether individual differences in geometry track an animal's behavior, and whether such a link is localized to particular brain regions or distributed across them.
To do this, we turned to public data released by the International Brain Lab \parencite{iblreproducibility}.
Two features of this dataset make it an ideal candidate for this analysis.
First, it consists of neural and behavioral measurements from a large cohort of animals performing the same binary choice task involving contrast discrimination (\textbf{Fig. }\ref{fig:ibl}a).
Second, a large number of neurons were measured in each animal spanning the same brain regions including secondary visual areas, hippocampus, and thalamus.

We included mice with more than n = 100 neurons and n = 40 behavioral trials (i.e. n = 20 trials per cross-validation fold), yielding a dataset of $K= 56$ subjects. As an input to our shape analyses, we computed 90-dimensional tuning curves (n = 9 contrast conditions and n = 10 time bins of 40 ms after the animal's response) for a total of n=24583 neurons. As done in previous analyses, we could detect differences across all recorded regions using Procrustes distances (\textbf{Fig.} \ref{fig:ibl}e), despite that regions recorded in this dataset have previously been shown to encode all task variables \parencite{iblreproducibility}.
To quantify the behavioral profiles on the task, we computed the probability of answering 'right' $P(\text{right})$ for each contrast.
These curves show considerable variability across subjects (\textbf{Fig.} \ref{fig:ibl}f), prompting us to ask whether there is a neural signature of this variability.

As before, we started by computing matrices of pairwise distances across subjects quantifying neural similarity using Procrustes, but also behavioral similarity using Pearson correlation between psychometrics curves (\textbf{Fig. }\ref{fig:ibl}g). A correlation analysis between behavioral and neural distances (\textbf{Fig.} \ref{fig:ibl}h) revealed a weak but significant relationship between these two measures (\textbf{Fig.} \ref{fig:ibl}d). Repeating this analysis using independent trials to calculate neural and behavioral distances between animals (thus avoiding spurious correlations due to shared noise) confirmed that the correlation values were consistently higher than the correlations between shuffled subjects (\textbf{Fig.} \ref{fig:ibl}i). This analysis confirmed that behavioral strategies are reflected in the encoding geometry. 

To pinpoint which region drove this relationship, we performed similar analyses conditioning on 5 different regions: visual area (VISa), CA1, dentate gyrus (DG), and lateral posterior (LP) and posterior (PO) thalamic nuclei. Specifically, we examined the correlation between neural and behavioral distances, as above, but using neural activity from each region separately, using mice with more than 40 trials in the analyzed region. We found significant differences for all regions but LP (\hl{XXX}). This analysis suggested that despite noticeable geometric differences between regions (\textbf{Fig.} \ref{fig:ibl}e), the relationship between behavioral strategies and encoding geometry was distributed across most regions. 

%Finally, we show once again (see also Figs. \ref{fig:synth} and \ref{fig:hdcells}) how aligning neural representations allows for zero-shot learning of psychometric curves using the two closest neighbors. Specifically, for every subject in the dataset we computed the behavioral correlation between this subject and the two nearest neighbor subjects in the neural shape space (closest, in \textbf{Fig.} \ref{fig:ibl}f). We then repeated this analysis, but now predicting each subject behavior using progressively distant subjects in neural shape space. We found that the behavioral correlation between the first (two closest neighbors) and the last (two furthest neighbors) was significantly different ($t=2.2, p=0.037$). Furthermore, we found a negative correlation between distance in neural shape space and behavioral distance (\textbf{Fig.} \ref{fig:ibl}f), showing that the behavioral correlation decreases smoothly for increasingly distant animals. All together, these analyses show how different behavioral strategies are reflected in different neural geometries, as quantified by shape metrics.

